% Self-contained mathematical core for the RAW parameter-substitution variant.
% Requires amsmath, amssymb, amsthm, theorem/lemma/corollary environments.
% Citation keys: templatefree, rawclausal, computationnormal, nativebiinterp.

\section{A raw calculus with parameter substitution}
\label{sec:rp-spec}

The principal simplification is a restriction on primitive equality
reasoning.  Proper-subterm substitution need only replace a free numerical
parameter.  Nonvariable sources are needed only for common-root equality
composition.  Arbitrary term-template transport is then derived by giving
its argument a fresh local name.  Unlike a presentation by computation
normal forms, this calculus retains ordinary raw syntax and ordinary
capture-avoiding substitution throughout.

The fresh-alias method is an existing alternative argument in
\cite{templatefree}, and clauses occur in \cite{rawclausal}.  The new
restriction is that arbitrary-pattern replacement can itself be removed:
the alias argument needs only variable substitution and root equality
composition.  The following rule basis, exact translations, and complete
checker comparison make that restriction concrete.

\subsection{Syntax, clauses, and ordinary readings}

Numerical terms are finite raw expressions
\[
 t::=x\mid0\mid S(t)\mid I(n,a;u.h),
\]
where only the displayed $u$ in $h$ is bound.  Formulas are
\[
 F::=(s=t)\mid D_{x,y}F.
\]
Both displayed coordinates of $D$ are bound.  All substitutions below
are ordinary capture-avoiding substitutions.  In particular,
$I(0,a;u.h)$ and $a$ are distinct terms, including after substitution.
There is no reduction or normalization while constructing a term or
substituting into a formula.

A reusable context is a finite list of clause occurrences
\[
 [A;y.B(y)],\qquad y\notin\operatorname{FV}(A).
\]
The ordinary interpretations are
\[
 (s=t)^\circ=(s=t),\qquad
 (D_{x,y}B)^\circ
   =\exists x\bigl(B^\circ(x,0)\to\forall y\,B^\circ(x,y)\bigr),
\]
\[
 [A;y.B]^\circ=A^\circ\to\forall y\,B^\circ(y).
\]
Equality has its actual numerical meaning.  Clauses are context objects,
not a third formula constructor.  Let $H=(0=0)$ and abbreviate an ordinary
assumption by the constant clause
\[
 \operatorname{ord}(F)=[H;v.F],\qquad v\notin\operatorname{FV}(F).
\]
The notation $\Gamma,F$ means $\Gamma,\operatorname{ord}(F)$.
Occurrences are retained even if their values coincide.  The designated
contradiction is $0=S0$.

\subsection{A fixed, restricted equality operation}

For raw terms $s,t,u$, define
\[
 \psi(s,t,u)=
 \begin{cases}
 u[t/s],&s\text{ is a free numerical parameter},\\
 t,&s\text{ is not a parameter and }u\equiv s,\\
 u,&s\text{ is not a parameter and }u\not\equiv s.
 \end{cases}
\]
The first case substitutes for a variable, with capture avoided, and does
not rescan an inserted target.  It is defined even if $t$ contains $s$.
For every source, $\psi(s,t,s)=t$.  For a nonvariable source, no proper
subterm is inspected; the last case merely reverses the minor equation
in the rule below.  Bound variables are not free numerical parameters.
The operation has no supplied equality template and no arbitrary-pattern
replacement pass.

Define also the partial one-step computation operation on an
iterator-headed term $b$:
\[
 \kappa(I(0,a;u.h))=a,\qquad
 \kappa(I(Sn,a;u.h))=h[I(n,a;u.h)/u].
\]
It is undefined for all other term shapes.  The recursive iterator in
the second result is left as raw syntax; $\kappa$ does not repeatedly
evaluate its result.

\subsection{Every primitive instruction}

The entire calculus $\mathcal P$ has the following five instruction
families:
\[
 \frac{}{\Gamma\vdash b=\kappa(b)}
 \quad\text{when }\kappa(b)\text{ is defined}.
 \tag{Compute}\label{rule:rp-compute}
\]
\[
 \frac{\Gamma\vdash s=t\qquad\Gamma\vdash u=r}
      {\Gamma\vdash r=\psi(s,t,u)}.
 \tag{Param}\label{rule:rp-param}
\]
\[
 \frac{\Gamma\vdash A}{\Gamma\vdash B(t)}
 \quad\text{if }[A;y.B(y)]\text{ is a selected occurrence in }\Gamma.
 \tag{Use}\label{rule:rp-use}
\]
\[
 \frac{\Gamma\vdash D_{x,y}B(x,y)\qquad
       \Gamma,[B(b,0);y.B(b,y)]\vdash F}
      {\Gamma\vdash F}.
 \tag{Open}\label{rule:rp-open}
\]
\[
 \frac{\Gamma,B(t,\eta)\vdash B(t,Sz)}
      {\Gamma\vdash D_{x,y}B(x,y)},\qquad
 \eta=\begin{cases}
 z,&B\text{ is a literal equation and }t\equiv0,\\
 0,&\text{otherwise}.
 \end{cases}
 \tag{Close}\label{rule:rp-close}
\]
In Open, $b$ is fresh for the retained context, conclusion, and all passive
parameters of $B$.  In Close, $z$ is fresh for the retained context,
witness, and all passive matrix parameters.  All terms, formulas, and
clauses must be well scoped.  The matrix-root and literal-zero tests are
tests on raw syntax.  These are part of the specification checked by the
program.

There is no separate reflexivity axiom, hypothesis rule, equality template,
first-class $Q$ constructor, numerical case rule, or induction schema.
Compute explicitly contains \emph{two} computation axiom schemas.  Calling
it one instruction family is an implementation grouping, not an assertion
that two mathematical equations became one equation.  Likewise the
finite cases of Param and Close are charged to the complete checker.

\section{The equality bootstrap and the fresh-alias theorem}
\label{sec:rp-bootstrap}

\begin{lemma}[Root equality algebra]
Common-left composition, reflexivity, symmetry, and transitivity are
admissible under any retained context:
\[
 s=t,\ s=r\vdash r=t,\qquad
 \vdash a=a,\qquad
 s=t\vdash t=s,\qquad
 s=t,\ t=r\vdash s=r.
\]
\end{lemma}
\begin{proof}
The first rule is a root instance of Param, since $\psi(s,t,s)=t$.
For reflexivity put $J(a)=I(0,a;u.0)$.  Compute proves $J(a)=a$;
using that equation as both premises of common-left composition gives
$a=a$.  For symmetry, use major $s=s$ and minor $s=t$.  For
transitivity, reverse $s=t$ and use major $t=r$ and minor $t=s$.
Every step uses only Compute and root equality composition.
\end{proof}

\begin{lemma}[Elementary proof operations]
Ordinary assumption use, weakening, injective fresh renaming, and cut
for ordinary formula assumptions are admissible.
\end{lemma}
\begin{proof}
Derived reflexivity proves $H$, so Use on $[H;v.F]$ proves $F$.
Weakening and injective fresh renaming follow by induction on a finite
derivation.  Such renaming preserves free-parameter status, exact root
matching, equation roots, and the literal zero.  For cut, distinguish the
occurrence $[H;v.F]$ being discharged.  Replace every Use selecting that
occurrence by the supplied proof of $F$, weakened into the local context.
At other instructions recurse through the premises, first freshening
eigenvariables away from the inserted proof.  This establishes cut
without presupposing arbitrary numerical substitution into derivations.
\end{proof}

\begin{lemma}[Fresh aliases]
\label{lem:rp-alias}
For a fixed raw term $a$ and $b$ fresh for $\Gamma$, $F$, and $a$,
\[
 \frac{\Gamma,b=a\vdash F}{\Gamma\vdash F}
 \tag{Alias}
\]
is admissible.  The local assumption is an ordinary formula assumption.
Only the two computation equations, root equality composition, the
logical instructions, and ordinary cut are needed to construct it.
\end{lemma}
\begin{proof}
Put $C(y,a,x)=I(y,a;u.x)$, with $u$ fresh, and let
\[
 B(x,y):=(C(y,a,x)=a).
\]
Close with witness $a$ proves $D_{x,y}B$.  Its successor branch
$C(Sz,a,a)=a$ is exactly the successor Compute instance with constant
update $a$.  The local assumption is unused, so this works whichever
mode Close selects, including when $a\equiv0$.

There is also a version using case-mode Close alone.  Take the raw
nonliteral witness $J(a)=I(0,a;v.0)$ instead.  Successor Compute gives
$C(Sz,a,J(a))=J(a)$, zero Compute gives $J(a)=a$, and root transitivity
joins them to the required branch.  This again ignores its local
assumption and always selects cases.  Thus the alias lemma does not
require equation progress, even though the shorter certificate can
use the direct witness $a$.

Open at fresh $b$.  Its retained clause is
\[
 [\,C(0,a,b)=a;\ y.C(y,a,b)=a\,].
\]
Zero Compute proves the antecedent.  Use at target $1=S0$ gives
$C(1,a,b)=a$.  Successor Compute gives $C(1,a,b)=b$.
Common-left composition, with the former as major and the latter as
minor, gives $b=a$.  Cut that ordinary assumption into the supplied
branch, after weakening it into the clause context.  Open discharges
$b$ and returns $F$ in the original context.
\end{proof}

This is the existing fresh-alias idea in a form adapted to the present
restricted equality interface.  Its significant feature for the
bootstrap is that it needs no proper-subterm substitution, term-template
transport, or equality induction.  It therefore cannot conceal the
equality principle that it is subsequently used to derive.

\begin{theorem}[Parameter substitution suffices for all templates]
\label{thm:rp-congruence}
For every fixed finite numerical term template $h(v)$,
\[
 \Gamma\vdash s=t
 \quad\Longrightarrow\quad
 \Gamma\vdash h(s)=h(t).
\]
Repeated, absent, and binder-nested occurrences of the displayed variable
are all permitted.
\end{theorem}
\begin{proof}
Work temporarily under the ordinary assumption $b=s$, with $b$ fresh
for the context, $h$, $s$, and $t$.  Transitivity gives $b=t$.
Derived reflexivity gives $h(b)=h(b)$.  Param with major $b=s$ and
this minor gives
\[
 h(b)=h(s),
\]
because substitution for the free variable $b$ marks exactly the
distinguished template occurrences.  The same construction with major
$b=t$ gives $h(b)=h(t)$.  Common-left composition joins these equations
to $h(s)=h(t)$.  Its endpoints contain no $b$, so
Lemma~\ref{lem:rp-alias} discharges the temporary alias.

Only the parameter case of Param is used for proper-subterm substitution.
The argument is not a pass replacing every occurrence of the possibly
compound term $s$ in $h(s)$: that would also replace passive occurrences
of $s$ already present in the template.  The untrusted proof compiler
builds $h(b)$, while each checker instruction receives only two proofs.
\end{proof}

\begin{corollary}[The complete equality interface]
Crossed one-term transport and transport in an arbitrary equation template
are admissible:
\[
 s=t,\ h(s)=r\vdash r=h(t),\qquad
 s=t,\ a(s)=b(s)\vdash a(t)=b(t).
\]
\end{corollary}
\begin{proof}
Join the congruence equation $h(s)=h(t)$ with the supplied minor to obtain
the crossed conclusion.  Apply crossed transport first to $a$ and then
to $b$ to obtain the second rule.
\end{proof}

In particular, the unrestricted one-pass replacement rule of
\cite{templatefree} is derived.  A root-first pass selects finitely many
pairwise disjoint source occurrences in a term $u$.  Replacing precisely
those occurrences by a fresh hole gives a finite template $h$ with
$h(s)=u$ and $h(t)$ equal to the pass's output.  The derived crossed
rule therefore gives the same conclusion.  This template extraction
can occur in an untrusted proof compiler; it is absent from the checker.

\section{Recovering the native logical rules}
\label{sec:rp-logic}

\begin{lemma}[Every case introduction]
\label{lem:rp-case}
For every matrix $B$ and witness $t$,
\[
 \frac{\Gamma,B(t,0)\vdash B(t,Sz)}
      {\Gamma\vdash D_{x,y}B(x,y)}
\]
is admissible with the usual freshness of $z$.
\end{lemma}
\begin{proof}
Only a literal equation $E$ at witness $0$ is not already Close's case
mode.  Work under a fresh alias $b=0$.  From the local $E(b,0)$,
equation-template transport gives $E(0,0)$.  Cut in the supplied branch,
obtaining $E(0,Sz)$.  Reverse the alias and transport back to
$E(b,Sz)$.  Close at the fresh variable $b$ is in case mode and gives
the required $D_{x,y}E(x,y)$.  Alias discharges $b$.
This uses only the already established atomic equality interface;
transport through a $D$ formula is not assumed.
\end{proof}

For fresh unused $w$, define the formula abbreviation
\[
 q_yB(y):=D_{w,y}B(y).
\]
Its three interfaces are derivable:
\begin{enumerate}
\item From $\Gamma,B(0)\vdash B(Sz)$ derive $\Gamma\vdash q_yB$.
      Close at unused witness $1$ is always in case mode.
\item From $q_yB$ and $B(0)$ derive $B(t)$.  Open the dummy witness,
      apply its clause by Use, and close Open.
\item From a retained clause $[B(0);y.B(y)]$ derive $q_yB$.  Under
      the first construction's local $B(0)$, use that clause at $Sz$.
\end{enumerate}
For a literal equation $E(y)$, Close at unused witness $0$ directly
gives $q_yE$ from $\Gamma,E(z)\vdash E(Sz)$.  With a supplied base
and the second interface, this is contextual equation induction.
The first two interfaces similarly give numerical cases for any formula.

Transport in every Eq/$D$ formula follows by external formula induction.
The equation case has been proved.  For $F(v)=D_{x,y}B(v,x,y)$, open
$F(s)$ at $b$ and introduce $F(t)$ at the same witness by cases.
Convert its target base $B(t,b,0)$ backwards using the induction
hypothesis and symmetry, use the retained source clause at $Sz$, and
convert the successor instance forwards.  This recursion uses only the
strictly smaller matrix.

\section{Exact consequence equivalence and iPRA}
\label{sec:rp-equivalence}

Let $\mathcal N$ denote the repository's raw native Eq/$Q$/$D$ calculus:
Hyp, the two iterator computations, one-term crossed transport, closure
elimination, paired elimination, case introduction, and direct
equation-progress introduction.  Its terms are exactly the raw terms
above, and its ordinary closure reading is
$Q_yB=B(0)\to\forall yB(y)$.

Define a formula translation, leaving numerical terms unchanged, by
\[
 \tau(s=t)=(s=t),\qquad
 \tau(D_{x,y}B)=D_{x,y}\tau B,\qquad
 \tau(Q_yB)=D_{w,y}\tau B,
\]
where each newly inserted witness coordinate is fresh and unused.
Let $j$ include Eq/$D$ formulas into native formulas.  The translation
commutes with ordinary capture-free numerical substitution and preserves
whether a formula's root is an equation or compound.

\begin{theorem}[All ordinary sequents are preserved and reflected]
For every finite native formula context and formula,
\[
 \Gamma\vdash_{\mathcal N}A
 \quad\Longleftrightarrow\quad
 \operatorname{ord}(\tau\Gamma)\vdash_{\mathcal P}\tau A.
\]
For Eq/$D$ formulas, $\tau jF=F$ literally up to bound renaming, and
for every native formula $A$,
\[
 A\ \dashv\vdash_{\mathcal N}\ j\tau A.
\]
\end{theorem}
\begin{proof}
Forwards, induct on the native derivation.  Hyp becomes a proof of $H$
and Use on the appropriate ordinary clause.  Computation becomes the
corresponding Compute case.  Transport uses
Theorem~\ref{thm:rp-congruence} and common-left composition.
Closure elimination uses the derived $q$ specialization.  Case
introduction uses Lemma~\ref{lem:rp-case}, and equation progress uses
Close at unused witness $0$.

For native paired elimination, the translated major is opened into its
clause.  The third closure interface derives the translation of the old
branch's $Q$ assumption from that clause.  Weaken the inductively
translated branch, cut in that ordinary assumption, and close Open.

Backwards, track each clause occurrence's origin.  An initial ordinary
clause, or an ordinary clause added in Close, is represented by the
corresponding native formula assumption.  A clause added in Open is
represented by the corresponding native $Q$ assumption.  Use on the
former is native Hyp, and Use on the latter is native closure elimination.
Compute is the matching native axiom.  Param is an instance of native
crossed transport: use template $u[v/s]$ if $s$ is a free parameter,
the identity template if $u\equiv s$ and $s$ is not a parameter, and
the constant template $u$ in the remaining case.  Here $v$ is chosen
fresh, so these substitutions avoid capture.

Open is native paired elimination.  A case-mode Close is native case
introduction.  A progress-mode Close at witness $0$ first uses native
equation progress to obtain $Q_yE(0,y)$; under $E(0,0)$ specialize it at
$Sz$ and introduce $D_{x,y}E(x,y)$ by native cases.  This proves reverse
preservation for all Eq/$D$ ordinary sequents.

The composite $\tau j$ is literal identity.  For the other composite,
the essential native equivalence is
\[
 Q_yB(y)\ \dashv\vdash\ D_{w,y}B(y),\qquad w\text{ unused}.
\]
Forwards, specialize the closure under its base and introduce the dummy
witness by cases.  Backwards, paired elimination opens exactly the
desired $Q$ assumption.  Uniform pointwise interderivability lifts
through $D$ by elimination and case introduction, converting the base
backwards and the successor instance forwards; it lifts through $Q$
in the same way with a dummy witness.  External structural induction
proves $A\dashv\vdash j\tau A$.  Ordinary cut with these equivalences
also gives the reverse implication for the displayed native sequent.
\end{proof}

The same proof represents an arbitrary externally supplied clause by
the native ordinary macro $A\to\forall yB(y)$ if a statement for arbitrary
clause contexts is desired.  The theorem already covers all ordinary
first-order consequences, not merely empty-context theorems.

\begin{corollary}[Equiconsistency and bi-interpretability]
The ordinary realization of $\mathcal P$ is equiconsistent and
bi-interpretable with full first-order intuitionistic primitive recursive
arithmetic.  Both numerical domains are whole, equality is actual
equality, and both numerical composite isomorphisms have graph $x=y$.
\end{corollary}
\begin{proof}
Compose the explicit interpretations and both formula round trips above
with the pinned repository's native/iPRA bi-interpretation
\cite{nativebiinterp}.  Numerical terms have not changed in the new
comparison, so its numerical maps are literal identities.  The inherited
bridge first recovers ordinary intuitionistic logic and the scalar
arithmetic theory with the iterator equations, coverage
\[
 \forall x\,(x=0\lor\exists y\,x=Sy),
\]
separation $0=S0\to\bot$, and contextual equation induction.  It then
constructs predecessor, reflection of equality, quantifier-free induction,
pairing, and scalar implementations of every fixed primitive-recursive
definition.  The resulting formula and term round trips, composed with
those above, establish the full first-order bi-interpretation.

The designated contradiction $0=S0$ is literally unchanged by $\tau$
and $j$, so finite contradiction proofs transfer in both directions.
The ordinary realization includes the stated falsity bridge; a theory
of positive equations alone would not exclude singleton structures.
\end{proof}

The new equality restriction and certificate compiler have written proofs
and executable checks, not a newly claimed Lean or Rocq formalization.
The inherited native/iPRA bridge is a separately formalized dependency.
No theory is asserted to prove its own consistency.

\subsection{Structural substitution and effective checking}

Arbitrary numerical substitution is not justified by copying instruction
tags: it can turn a Param source from a variable into a compound term,
and can turn a nonzero-spelled Close witness into the literal zero.
Substitution admissibility for ordinary sequents instead follows from the
exact comparison with the established native calculus.  The bootstrap
above used only weakening, fresh renaming, and ordinary-assumption cut,
which were proved directly.

The raw certificate checker uses finite structural traversals,
capture-free substitutions, literal comparisons, and one-step computation.
It does not call a solver or a uniform evaluator for primitive-recursive
programs.  Under standard finite-tree encodings its verification relation
is primitive recursive.  This does not imply a useful polynomial bound:
proof trees and substituted terms can still grow substantially.
The smaller source comparison must count parsing, scope checks, clause
formation and opening, the two Compute cases, all Param cases, and every
eigenvariable check.  A five-instruction presentation by itself is not
a theorem of global minimality.

